% ===== 原文 pp.67–70（PDF p.77–80）及印刷 p.71 顶部（PDF p.81，仍属 §6.4.1.2）=====
% [原文 p.67]
更一般地，若共轭类不是闭的，则需考虑其闭包的余法丛。

容易分析这些边界条件的边界线算子。把 3d 理论在带有特殊地位边界条件的区间上约化，我们得到 B-型 2d 规范理论耦合到以 $G_\C$ 为靶空间的 B-模型。规范群通过共轭作用作用在 $G_\C$ 上。根据 A.3，相应的膜范畴等价于 $D^b_{G_\C}(Coh(G_\C))$。幺半结构无法由 2d 考虑推出，但与通常 RW 模型中同样的分析表明它由导出张量积给出。

% [Tutor Note: p.67 余法丛 $N^*\cC$ 的定维核对：$T^*G_\C$ 复维为 $2n$（$n=\dim_\C G_\C$）；复子流形 $\cC$ 的余法丛复维为 $\dim_\C\cC+(n-\dim_\C\cC)=n$，恰为全空间复维的一半，故必为复拉格朗日——这解释了为何"取余法丛"这一构造对任何复子流形都自动满足拉格朗日条件，$G_\C$-不变性只需 $\cC$ 是共轭类。]

在 Gukov-Witten 情形，我们需要固定 $G_\C$ 中一个半单复共轭类 $\cC$。用 $N^*\cC$ 记它在 $T^*G_\C$ 中的余法丛的全空间。具体地，它是点对 $(g,\sigma)$ 构成的空间，其中 $g\in\cC$，而 $\sigma\in\fg_\C$ 对任意与 $\cC$ 在 $g$ 处相切的 $\delta g$ 都满足 $\Tr\sigma g^{-1}\delta g=0$。纤维坐标 $\sigma$ 的上同调次数为 2；为表明这一点，我们把相应的分次复流形记作 $N^*[2]\cC$。Gukov-Witten 情形下在区间上的约化给出 B-型 2d 规范理论耦合到一个 B-模型，其靶空间是 $N^*[2]\cC$。它的膜范畴是 $D^b_{G_\C}(Coh(N^*[2]\cC))$。幺半结构由导出张量积给出。

% [Tutor Note: p.67 余法条件 $\Tr\sigma g^{-1}\delta g=0$ 的含义：把 $\sigma\in\fg_\C$ 与 $g^{-1}\delta g\in T_g\cC\subset\fg_\C$ 用不变双线性型配对，条件即"$\sigma$ 与共轭类的切空间正交"，也就是 $\sigma$ 落在余法空间里；改用 $\fg_\C^*$ 的语言即 $\sigma$ 在 $T_g\cC$ 上为零。]

\subsection{体线算子}

考虑与 $G_\C$ 中平凡共轭类（即单位元）对应的 Gukov-Witten 面算子的特殊情形是有意思的。这就是平凡面算子，因此此时 3d 边界线算子的范畴可以等同于 4d TFT 中的体线算子范畴。单位元的余法丛就是复化李代数 $\fg_\C$ 的对偶；群 $G_\C$ 通过伴随表示作用在它上面。因此 4d 体线算子的范畴等价于 $D^b_{G_\C}(Coh(\fg^*_\C[2]))$。换句话说，它是分次代数 $\oplus_p\operatorname{Sym}^p\fg$ 的 $G_\C$-等变导出范畴，其中第 $p$ 个分量处于上同调次数 $2p$。

% [Tutor Note: p.67 从 $\cC=\{\id\}$ 到 $\fg_\C^*$ 的核对：单位元的共轭类只是一个点，其余法丛即该点处的余切空间 $T^*_{\id}G_\C\cong\fg_\C^*$；由伴随不变性，$G_\C$ 在 $\fg_\C^*$ 上的作用正是余伴随（伴随的对偶）表示。$\fg_\C^*$ 的坐标环是 $\operatorname{Sym}^\bullet\fg$，而 "[2]" 表示每升高一次 $\operatorname{Sym}$ 的次数就升高上同调次数 2，故第 $p$ 个分量在次数 $2p$。]

鉴于这一结果，考虑位于两个 Wilson 圈交接处的局部算子是有意思的，这两个 Wilson 圈分别取 $G$ 的表示 $V_1$ 与 $V_2$。范畴 $D^b_{G_\C}(Coh(\fg^*_\C[2]))$ 中相应的对象是 $\mathfrak{A}=\oplus_p\operatorname{Sym}^p\fg[2]$ 上的自由模，形如 $V_1\otimes_\C\mathfrak{A}$ 与 $V_2\otimes\mathfrak{A}$，带显然的 $G_\C$ 作用。它们之间的态射空间是无穷维分次表示
\[
V_1^*\otimes V_2\otimes\mathfrak{A}
\]
中 $G_\C$-不变量的空间。确实，两条 Wilson 线的 BRST 不变且规范不变的交接点，应当是表示 $V_1^*\otimes V_2$ 中由取值于 $\fg_\C$ 的复标量 $\sigma$ 构造出来的算子。这类算子在鬼数 $2p$ 处的空间是 $\Hom_G(\operatorname{Sym}^p\fg, V_1^*\otimes V_2)$，其中 $\Hom_G$ 表示 $G$ 的表示范畴中的态射空间。对所有 $p$ 求和就得到上面的答案。

% [Tutor Note: p.67 态射空间两种算法的等价性：有限维表示下 $(V_1^*\otimes V_2\otimes\operatorname{Sym}^p\fg)^G\cong\Hom_G((\operatorname{Sym}^p\fg)^*, V_1^*\otimes V_2)$，而 $\fg$ 上有 $G$-不变配对使 $\operatorname{Sym}^p\fg$ 自对偶，故等于原文的 $\Hom_G(\operatorname{Sym}^p\fg, V_1^*\otimes V_2)$。鬼数 $2p$ 与 $\operatorname{Sym}^p\fg$ 的对应来自 $\sigma$ 的次数 2。]

% [原文 p.68]
\subsection{更一般的面算子}

与阿贝尔情形一样，上面的例子并未穷尽面算子 2-范畴中的对象。例如，[47] 中考虑了更复杂的面算子，它们涉及复联络 $\cA=A+\ii\phi$ 的高阶极点。类比 Rozansky-Witten 模型，我们提出：$t=\ii$ 时最一般的面算子（等价地，3d 理论中最一般的边界条件）可以定义为有特殊地位边界条件 $\mathbb{X}_0$ 的边界线算子幺半范畴上的一个模范畴。如上所述，这个幺半范畴是 $D^b_{G_\C}(Coh(G_\C))$。具体地，这意味着最一般的面算子可以这样得到：把一族 2d TFT 纤维化在 $G_\C$ 上，使得底空间上的 $G_\C$ 作用（由共轭给出）提升为整个族上的 $G_\C$ 作用。例如，可以取一个复流形 $X$，它是 $G_\C$ 上的纤维化，纤维是 Calabi-Yau 流形，并且给定了底空间上 $G_\C$ 作用（由共轭给出）到全空间上 $G_\C$ 作用的一个提升。

给定任意面算子 $\mathbb{X}$，可以通过考察位于 $\mathbb{X}$ 与有特殊地位面算子 $\mathbb{X}_0$ 交接处的线算子范畴，构造 $D^b_{G_\C}(Coh(G_\C))$ 上的一个模范畴。这个范畴是把 3d TFT 在区间上紧化所得的 2d TFT 中的膜范畴，其中区间两端的边界条件分别由 $\mathbb{X}$ 与 $\mathbb{X}_0$ 给出。等价地，可以把 4d TFT 在带两个穿刺点的 2-球面上紧化，并在两个穿刺点处插入面算子 $\mathbb{X}$ 与 $\mathbb{X}_0$。

% [Tutor Note: p.68 第二种等价描述（带两个穿刺点的 2-球面）的核对：把区间 $I$ 两端的 3d 边界条件换成穿刺点，用的是"4d 理论在 2-球面上紧化、面算子沿穿刺点截断"这一标准对应——面算子在此恰好给出 2d 理论中的线算子，而"区间 + 两端边界条件"是它的 3d 对偶图像。两边得到的都是同一个模范畴（$\mathbb{X}$ 与 $\mathbb{X}_0$ 之间的缺陷线范畴）。]

例如，若按 [47] 的方式，用复联络 $\cA$ 的一个给定奇点来定义面算子，并注意到有特殊地位的面算子是靠允许 $\cA$ 的和乐自由取值来定义的，就会看到有效 2d TFT 的真空空间是带给定奇点的穿孔圆盘上联络的模空间。把这个模空间记作 $\cM$。如果在 $\cM$ 的定义中再模去那些在圆盘边界上某个选定点处化为恒等的规范变换群，则 $\cM$ 被 $G_\C$ 作用，并纤维化在 $G_\C$ 上（后者即 $\cA$ 沿圆盘边界的和乐）。有效 2d TFT 看来是以 $\cM$ 为靶空间、耦合到规范群为 $G_\C$ 的 B-型规范理论的 B-模型，这是颇为合理的。它的膜范畴是 $D^b_{G_\C}(Coh(G_\C))$ 上的一个模范畴。

% [Tutor Note: p.68 原文有笔误 "consider a complex manifolds $X$"，应为 "a complex manifold $X$"（单数）。内容按复数流形译出，不影响数学。]

\section{\texorpdfstring{$t=1$}{t=1} 且 \texorpdfstring{$G=U(1)$}{G=U(1)} 时的面算子}

在 §5.4 中我们发现，对 $t=1$，GL-扭转理论紧化为以 $G$ 为靶空间的带规范场 3d A-模型 TFT。物质扇区（以 $G$ 为靶空间的 3d A-模型，见 §B.1）的边界条件分析已在 [6] 中完成。然而，完整带规范场 3d A-模型的边界条件此前尚未被分析过。本节中我们限于 $G=U(1)$，此时规范扇区与物质扇区干净地分离。

% [Tutor Note: p.68 为何 $G=U(1)$ 时两个扇区可以分开：此时 3d A-模型的靶空间 $G=U(1)=S^1$ 只有一维，物质扇区的场（取值于 $G$ 的映射）与规范场仅通过 $A_4$（约化后的周期标量）耦合，规范扇区可以独立地取 Dirichlet 或 Neumann 条件（§6.4.1），物质扇区另行处理（§6.4.2）。核对时注意这里分析的是 3d 约化后的理论，不是 4d 原理论。]

% [Tutor Note: p.68 "以 $G$ 为靶空间"的来源核对：$t=1$ 扭转后，$\cN=4$ 的标量场组合成取值于 $G$ 的映射；$G=U(1)$ 时靶空间就是圆周 $S^1$，其坐标为角变量。这一点与 p.71 出现的"周期标量 $A_4$"相呼应——周期标量正是靶空间 $S^1$ 的坐标。若把 §5.4 的约化谱与此处对照，靶空间维数应与标量场个数一致。]

% [原文 p.69]
\subsection{规范扇区中的边界条件}

规范扇区是 Donaldson-Witten 4d TFT [45] 到 3d 的维数约化；在附录 B 中我们把它称为 A-型 3d 规范理论。然而仅此一点并不能教给我们很多东西，因为这个理论中的边界条件此前没有被讨论过。在不添加边界自由度的情况下，仅有的选择是规范场的 Dirichlet 与 Neumann 边界条件，而 BRST 不变性会固定所有其他场的条件。

% [Tutor Note: p.69 "不添加边界自由度时只有 D/N 两种"的依据：3d 规范场 $A$ 在边界上的法向分量为 $A_3$，其共轭动量是 $\star F$ 的相应分量；仿照电动力学，要么固定 $A_i$（Dirichlet）、要么固定其共轭动量（Neumann），二者不能同时固定，否则边界条件不相容。BRST 不变性（超对称的残余）接着唯一地确定费米子与 $\phi_4$ 的条件。]

% [Tutor Note: p.69 "A-型 3d 规范理论"的来历：它是 Donaldson–Witten 4d TFT [45]（$\cN=2$ 超对称杨-米尔斯经特定扭转）维数约化到 3d 的结果；其超荷满足 $\delta_Q^2=2\ii\delta_g(\sigma)$（p.69 即用到），故标量 $\sigma$ 参数化了超荷平方所生成的规范变换。与 §B.4 的"B-型 3d 规范理论"相对照，两者的差别在于所选拓扑超荷不同，这决定了边界条件与可观测量的结构。]

\subsubsection{Dirichlet 条件}

我们从 Dirichlet 条件开始，它说 $A$ 在边界上的限制是平凡的。由于 $\delta_Q^2=2\ii\delta_g(\sigma)$，这一条件只有在 $\sigma$ 也于边界上为零时才有意义。BRST 不变性随后要求 $\eta+\tilde\eta$ 与 1-形式 $\psi+\tilde\psi$ 的限制为零。费米子运动方程随后要求 $\chi$ 的限制为零，而 BRST 不变性意味着 $\phi_4$ 必须满足 Neumann 条件 $\partial_3\phi_4=0$，其中我们假定边界由 $x^3=0$ 给出。

% [Tutor Note: p.69 为何 Dirichlet 条件要求 $\sigma|_{\partial}=0$：把 $\delta_Q^2=2\ii\delta_g(\sigma)$ 作用在被 Dirichlet 条件固定的 $A$ 上，结果必须仍为零；而 $\delta_g(\sigma)$ 作用会把 $A|_{\partial}=0$ 变成 $-\dd_A\sigma|_{\partial}$，故须 $\sigma$ 在边界上为零（至少为常数，再结合 $\phi_1,\phi_2$ 的条件得零）。这是原文"makes sense only if"处省略的论证。]

Dirichlet 边界条件具有这样的性质：它没有非平凡的局部 BRST 不变边界可观测量。确实，唯一非零的 BRST 不变 0-形式是 $\psi_4-\tilde\psi_4$，但它是 BRST-正合的。为分析边界线算子，我们用维数约化的技巧，把 3d 理论在带 Dirichlet 边界条件的区间上紧化。有效 2d 理论中仅有的玻色场是 $\phi_4$ 的常数模，以及 $A$ 沿由 $x^3$ 参数化的区间的和乐。也就是说，玻色场是一个实标量和一个周期实标量。因此有效 2d TFT 是以 $\R\times S^1$ 为靶空间的 $\sigma$-模型。事实上，它可以看作以 $T^*S^1$ 为靶空间的 A-模型。看出这一点的最容易的方法是注意到：3d 理论的路径积分局域化在 Bogomolny 方程
\[
F+\star\dd\phi_4=0
\]
的解所给出的构型上。在除 $A_3$ 与 $\phi_4$ 之外的所有场都置零、并假定它们与 $x^3$ 无关之后，这个方程变成
\[
\dd A_3+\star\dd\phi_4=0,
\]
其中 $\star$ 是 2d 的 Hodge 星算子。这是一个椭圆方程，只要宣布 $A_3+\ii\phi_4$ 是靶空间上的复坐标，它就可以解释成全纯瞬子方程。由于 4d 理论的作用量是 BRST-正合的，2d 模型的作用量也是。这与一个众所周知的事实相符：若靶空间上的辛形式是恰当的，则 A-模型的作用量是 BRST-正合的。

% [Tutor Note: p.69 从 $\dd A_3+\star\dd\phi_4=0$ 到"全纯瞬子方程"的补算：设 $J$ 是底流形上的复结构，2d 中在 1-形式上有 $\star\omega=\omega\circ J$（差一个符号约定），则该式等价于 $(\dd A_3+\ii\dd\phi_4)\circ J=-\ii(\dd A_3+\ii\dd\phi_4)$，即 $A_3+\ii\phi_4$ 作为靶空间 $\C$ 中的坐标满足 Cauchy--Riemann 方程 $\dbar(A_3+\ii\phi_4)=0$，也就是全纯映射。这正是"只要宣布 $A_3+\ii\phi_4$ 为复坐标"的含义。]

% [Tutor Note: p.69–70 两句的衔接核对：$T^*S^1\cong S^1\times\R$ 的辛形式 $\dd(p\,\dd\theta)$ 整体恰当，故 A-模型作用量 BRST-正合；而其镜像（SYZ 对偶）为 $\C^*$，这给出 p.70 所用 $Fuk(T^*S^1)\cong D^b(Coh(\C^*))$。两处说法必须一致，核对时以该镜像等价为准。]

Dirichlet 边界上的线算子范畴因此是 $T^*S^1$ 的 Floer--Fukaya 范畴，其最简单的对象是带平坦联络的酉向量丛所配的拉格朗日子流形。

% [原文 p.70]
由于这个范畴是某个 2-范畴中一个对象自同态范畴，它必定具有幺半结构，而这一结构从纯粹的 2d 视角是看不见的。事实上，我们并不指望一般辛流形的 Floer--Fukaya 范畴有自然的幺半结构。我们将在下文论证：该幺半结构是由镜像对称诱导的，后者建立了 $T^*S^1$ 的 Floer--Fukaya 范畴与 $D^b(Coh(\C^*))$ 的等价，并把后一范畴上的幺半结构搬到前者。眼下我们只注意：底空间 $S^1$ 有一个特殊点，对应于 $A$ 在区间上的平凡和乐。该点上的纤维是 $T^*S^1$ 中的一个拉格朗日子流形，它是幺半结构下的单位对象。这个特殊点让我们可以把 $S^1$ 与群流形 $U(1)$ 等同起来。

% [Tutor Note: p.70 原文 "allows us to identity $S^1$ with the group manifold $U(1)$" 中 "identity" 应为 "identify"（识别）。判断依据：此处是"等同"而非"恒等"之义，且 "identity object" 在同句中已用过，属形近词之误。译文按 "identify" 处理。]

\subsubsection{Neumann 条件}

现在考虑 3d 规范场 $A$ 的 Neumann 条件。这意味着规范对称性在边界上未被破缺，且 1-形式 $\star F$ 的限制为零。于是 Bogomolny 方程要求 $\phi_4$ 满足 Dirichlet 边界条件 $\phi_4=a=\mathrm{const}$，而由 BRST 不变性，$\psi_4-\tilde\psi_4$ 必须在 $x^3=0$ 处为零。费米子运动方程随后意味着 $\psi_3+\tilde\psi_3$ 也为零；又由于 $\delta_Q(\psi_3+\tilde\psi_3)=2\partial_3\sigma$，场 $\sigma$ 满足 Neumann 条件。最后，为使费米子边界条件相容，1-形式 $\star\chi$ 在边界上的限制必须为零。确实，若把 $x^1$ 视为时间方向，则 $(\star\chi)_2$ 与 $\psi_3+\tilde\psi_3$ 是正则共轭的，所以若其中一个为零，另一个也应当为零。类似地，若把 $x^2$ 视为时间，则 $(\star\chi)_1$ 与 $\psi_3+\tilde\psi_3$ 正则共轭，因而也必须为零。

% [Tutor Note: p.70 由 Neumann + Bogomolny 推出 $\phi_4=a$ 的补算：把 $F+\star\dd\phi_4=0$ 用 $\star$ 作用一次并限制到边界，得 $\star F\pm\dd\phi_4=0$；由 Neumann 条件 $\star F|_{\partial}=0$ 即得 $\dd\phi_4|_{\partial}=0$，故 $\phi_4$ 在边界上为常数 $a$。这是原文"the Bogomolny equation requires"一句跳过的论证。]

在 Neumann 情形，边界上 BRST 不变局部可观测量的空间由场 $\sigma$ 的幂张成。要确定边界线算子的范畴，需把 3d 规范理论在带 Neumann 边界条件的区间上约化。有效 2d 理论的玻色场是 2d 规范场与标量 $\sigma$ 的常数模，费米子场是 0-形式 $\eta+\tilde\eta$、1-形式 $\psi+\tilde\psi$ 与 2-形式 $\chi$。它们的 BRST 变换为
\[
\begin{aligned}
\delta_QA&=\ii(\psi+\tilde\psi), &
\delta_Q\sigma&=0, &
\delta_Q\bar\sigma&=\ii(\eta+\tilde\eta),\\
\delta_Q(\eta+\tilde\eta)&=0, &
\delta_Q(\psi+\tilde\psi)&=2\dd\sigma, &
\delta_Q\chi&=F .
\end{aligned}
\]
这个 2d TFT 可以由通常的 $\cN=(2,2)$、$d=2$ 超对称规范理论通过
% [原文 p.71]
一个利用 $U(1)_V$ R-对称性的扭转而得到。由于这与构造 A-型 $\sigma$-模型所用的 R-对称性相同，我们不妨把这个 TFT 称为 A-型 2d 规范理论。据我们所知，此前文献中尚未分析过它的边界条件。§A.1 表明，它的膜范畴等价于分次直线 $\C[2]$ 上凝聚层的有界导出范畴。\footnote{这个范畴等价于单点上层的 $U(1)$-等变可构造导出范畴 [28]。}再者，这个范畴的 3d 来源意味着它必定具有幺半结构。这里的幺半结构由凝聚层复形的通常导出张量积给出。$\C[2]$ 上的平凡线丛是单位对象。从 3d 视角看，它对应于边界上"不可见"的线算子。

% [Tutor Note: p.70 BRST 变换表的自洽性核对：$\delta_Q^2\sigma=0$、$\delta_Q^2(\eta+\tilde\eta)=0$ 自动成立；$\delta_Q^2(\psi+\tilde\psi)=2\delta_Q\dd\sigma=0$；$\delta_Q^2\chi=\delta_QF$ 应与 $2\ii\delta_g(\sigma)$ 相容（在 $G=U(1)$ 时 $\delta_g(\sigma)$ 平凡），故表是自洽的。此表与 §A.1 中 A-型 2d 规范理论的变换一致，这也是下文把该 TFT 称为 A-型 2d 规范理论的依据。]

% [Tutor Note: p.70 参数空间核对（见 p.71）：Neumann 条件在量子层面除边界值 $a\in\R$ 外还有 $\theta\in\R/2\pi\Z$ 的拓扑项系数，合起来 $\R\times S^1\simeq\C^*$；这与 §6.2.3 中 Gukov–Witten 参数 $h_0\in\C^*$ 的取值域一致（$h_0=\exp(-2\pi\tau)$），是 $t=1$ 与 $t=\ii$ 两面描述可以拼合的迹象。]

% ---------------------------------------------------------------
% 以下为印刷 p.71 顶部、仍属 §6.4.1.2 的正文（至 §6.4.2 标题之前）。
% 说明：§6.4.1.2 在原文中并未于 p.70 结束，而是延续到 p.71 顶部；
% 为免跨页句子被截断、并使 §6.4.1.2 完整，此处一并译出，止于 §6.4.2 之前。
% ---------------------------------------------------------------
% [原文 p.71]

如上所述，Neumann 条件依赖于一个实参数 $a$，即标量 $\phi_4$ 的边界值。在量子层面还有另一个参数，取值于 $\R/2\pi\Z$。它作为边界作用量中一个拓扑项的系数出现：
\[
\theta\int_{x^3=0}\frac{F}{2\pi}.
\]
于是总体上，规范扇区中的 Neumann 条件具有参数空间 $\R\times S^1\simeq\C^*$。

% ===== 本块结束（原文 §6.4.2 起属 tr-06 的块，此处不写） =====
